%----------------------------------------------------------------------------------------
% SECTION TITLE
%----------------------------------------------------------------------------------------

\cvsection{انتشارات و اختراعات}

%----------------------------------------------------------------------------------------
% SECTION CONTENT
%----------------------------------------------------------------------------------------

\begin{cventries}

\cventry
{ثبت اختراع آمریکا شماره 11,429,712}
{Systems and methods for dynamic passphrases}
{USPTO}
{اوت ۲۰۲۲}
{
\begin{cvitems}
\item {این اختراع سامانه‌ای امن و کارآمد برای تولید و مدیریت گذرواژه‌های پویا ارائه می‌کند که می‌تواند امنیت سامانه‌های احراز هویت و حفاظت از حساب‌ها و داده‌های کاربران را در برابر دسترسی غیرمجاز بهبود دهد.}
\item {\textbf{اثرگذاری:} افزایش امنیت و بهبود تجربه کاربری در سامانه‌های احراز هویت با حذف نیاز به حفظ گذرواژه‌ها یا PINهای پیچیده.}
\end{cvitems}
}

\cventry
{پذیرفته‌شده — چهارمین کارگاه بین‌المللی تفسیرپذیری هوش ماشین در پردازش تصاویر پزشکی، ۲۰۲۱}
{Soft Attention Improves Skin Cancer Classification Performance}
{استراسبورگ، فرانسه}
{سپتامبر ۲۰۲۱}
{
\begin{cvitems}
\item {ارائه سازوکار جدیدی برای توجه نرم که می‌تواند در سایر مسائل طبقه‌بندی تصویر نیز به کار رود و زمینه پژوهش‌های بیشتر در استفاده از سازوکارهای توجه در یادگیری عمیق را فراهم می‌کند.}
\item {\textbf{اثرگذاری:} افزایش دقت طبقه‌بندی سرطان پوست با سازوکار توجه نرم پیشنهادی، با پتانسیل کاربرد در ابزارهای تشخیصی و سامانه‌های تشخیص به کمک رایانه.}
\end{cvitems}
}

\cventry
{پذیرفته‌شده — IPMI 2021}
{Improving Joint Learning of Chest X-Ray and Radiology Report by Word Region Alignment}
{استراسبورگ، فرانسه}
{سپتامبر ۲۰۲۱}
{
\begin{cvitems}
\item {ارائه شبکه‌ای برای یادگیری مشترک بازنمایی تصویر و متن با هدف پیش‌آموزش روی تصاویر رادیوگرافی قفسه سینه و گزارش‌های رادیولوژی.}
\item {\textbf{اثرگذاری:} رمزگذار تنظیم‌شده با مقداردهی اولیه از بازنمایی بالینی پیش‌آموزش‌دیده پیشنهادی، عملکرد طبقه‌بندی بهتری را در زمان آموزش کوتاه‌تر نسبت به طبقه‌بندی چندبرچسبی نظارت‌شده به دست آورد.}
\end{cvitems}
}

\cventry
{ثبت اختراع آمریکا شماره 11,539,525}
{Systems and methods for secure tokenized credentials}
{USPTO}
{ژانویه ۲۰۲۰}
{
\begin{cvitems}
\item {ارائه روشی برای توکن‌سازی اطلاعات احراز هویت کاربر، مانند نام کاربری و گذرواژه، و تولید توکنی که برای احراز هویت بدون افشای اطلاعات حساس اصلی استفاده می‌شود. این سامانه همچنین شامل خزانه امن برای ذخیره و مدیریت توکن‌هاست.}
\item {\textbf{اثرگذاری:} افزایش امنیت سامانه‌های احراز هویت آنلاین با کاهش خطر نشت داده و سرقت هویت و فراهم‌کردن کاربردهای تجاری برای سازمان‌هایی مانند بانک‌ها و سامانه‌های تجارت الکترونیک.}
\end{cvitems}
}

\cventry
{پذیرفته‌شده — هفدهمین کنفرانس بین‌المللی پردازش زبان طبیعی (ICON 2020)}
{Self-Supervised Claim Identification for Automated Fact Checking}
{IIT Patna، هند}
{دسامبر ۲۰۲۰}
{
\begin{cvitems}
\item {ارائه رویکردی جدید، مبتنی بر توجه و خودنظارتی، برای شناسایی جملات دارای قابلیت طرح ادعا در یک مقاله خبری جعلی؛ گامی مهم در راستی‌آزمایی خودکار اخبار.}
\item {\textbf{اثرگذاری:} راهبرد یادگیری خودنظارتی مبتنی بر توجه، در تشخیص جملات کلیدی و درک محتوای اصلی مقاله خبری عملکرد بهتری ارائه می‌دهد.}
\end{cvitems}
}

\cventry
{پذیرفته‌شده — ICFHR 2020 | \textbf{DOI:} 10.1109/ICFHR2020.2020.00074}
{\href{https://arxiv.org/abs/2009.04532}{Attention based Writer Independent Handwriting Verification}}
{دورتموند، آلمان}
{ژوئیه ۲۰۲۰}
{
\begin{cvitems}
\item {ارائه مدل جدید مبتنی بر توجه برای محاسبه احتمال تعلق نمونه‌های دست‌نوشته اسکن‌شده به یک فرد واحد.}
\item {\textbf{اثرگذاری:} دستیابی به نتایج SOTA در تطبیق نمونه‌ها با استفاده از راهبرد توجه متقاطع و توجه نرم برای شناسایی نواحی مهم در تصاویر دست‌نوشته؛ این روش قابلیت استفاده در احراز هویت چهره و اثر انگشت را نیز دارد.}
\end{cvitems}
}

\cventry
{پذیرفته‌شده — کارگاه Interpretable and Explainable Machine Vision (BMVC 2019) | \textbf{arXiv:} 1909.02548v1}
{\href{https://arxiv.org/pdf/1909.02548.pdf}{Explanation based Handwriting Verification}}
{کاردیف، بریتانیا}
{سپتامبر ۲۰۱۹}
{
\begin{cvitems}
\item {ارائه معماری مؤثر Hybrid Deep Learning (HDL) برای تعیین احتمال اینکه یک واژه دست‌نویس مورد بررسی توسط نویسنده شناخته‌شده نوشته شده باشد.}
\item {\textbf{اثرگذاری:} انتشار یک مجموعه‌داده معیار برای پژوهش در زمینه احراز هویت دست‌خط مبتنی بر توضیح‌پذیری.}
\end{cvitems}
}

\cventry
{پذیرفته‌شده — ICFHR 2018 | \textbf{DOI:} 10.1109/ICFHR-2018.2018.00040}
{\href{https://www.researchgate.net/profile/Mohammad_Abuzar_Shaikh/publication/329564315_Writer_Verification_using_CNN_Feature_Extraction/links/5c658214a6fdccb608c3a432/Writer-Verification-using-CNN-Feature-Extraction.pdf}{Writer Verification Using CNN Feature Extraction}}
{نیاگارا، نیویورک، آمریکا}
{اوت ۲۰۱۸}
{
\begin{cvitems}
\item {ارائه روش یادگیری سرتاسری مبتنی بر ویژگی‌های آماری استخراج‌شده در سطح مجموعه نمونه‌ها برای حل مسئله احراز هویت نویسنده؛ یعنی تعیین یکسان‌بودن دو منبع دست‌خط بر اساس نمونه‌های هر دو منبع.}
\item {\textbf{اثرگذاری:} اثبات تجربی عملکرد بهتر ویژگی‌های آماری پارامتریک نسبت به آزمون‌های استاندارد مقایسه توزیع.}
\end{cvitems}
}

\cventry
{پذیرفته‌شده — ICFHR 2018 | \textbf{DOI:} 10.1109/ICFHR-2018.2018.00041}
{\href{https://arxiv.org/pdf/1812.02621.pdf}{Hybrid Feature Learning for Handwriting Verification}}
{نیاگارا، نیویورک، آمریکا}
{اوت ۲۰۱۸}
{
\begin{cvitems}
\item {ارائه معماری مؤثر Hybrid Deep Learning (HDL) برای تعیین احتمال نوشته‌شدن یک واژه دست‌نویس توسط نویسنده شناخته‌شده.}
\item {\textbf{اثرگذاری:} اثبات تجربی عملکرد بهتر ترکیب ویژگی‌های SIFT دستی با بازنمایی‌های یادگیری عمیق نسبت به مدل‌های صرفاً مبتنی بر ویژگی‌های دستی.}
\end{cvitems}
}

\end{cventries}
